Neural computer development has accelerated along three converging threads.

\textbf{Video world models and generators.}
High‑capacity video systems (e.g., Sora, Veo, Movie Gen) show that learned dynamics can track complex scenes over dozens of seconds.
This suggests that a video‑generation backbone can serve as the basis for interface modeling when properly conditioned on inputs and actions.

\textbf{Operating‑system and desktop modeling.}
Neural‑OS and related efforts highlight desktop‑scale simulation with recurrent and latent diffusion components, pushing toward OS‑level behavior.
These systems point to feasibility but often rely on specialized pipelines; our approach favors a mainstream, unified backbone to maximize transfer across interfaces.

\textbf{Agent systems.}
DOM‑centric agents succeed on structured web UIs with strong instrumentation and evaluation.
We treat structural signals as auxiliaries rather than the core system.
A pixel/latent NC generalizes across web and non‑web interfaces while still benefiting from structure when available.

Open challenges remain central to NCs: long‑term memory, robustness under low consistency (ambiguity, occlusion, nondeterminism), variable‑rate scheduling, and cross‑application generalization.
Our implementation targets these gaps with a pixel‑grounded NC and functional evaluation signals tied to on‑screen reality.
